\chapter{Introduction}
\label{chap-introduction}

\section{Context}

The course \emph{Initiation aux interfaces graphiques} at Sup Galilée~\cite{toulouse2023}
ends with a project: giving a Battleship game a graphical interface with the GTK library.
The course is not just about learning windows and buttons. It argues for a way of organising
code, \emph{modular programming}, which its material returns to from one session to the next
and which can be summed up as follows:

\begin{itemize}
    \item for each concept, a header file that gives its definition and how to use it, and a
          source file that implements the functions that manipulate it;
    \item using a concept is independent of how it is implemented;
    \item one program can be given several interfaces --- console, window --- and that is
          then the only part that has to change;
    \item every function written is tested straight away.
\end{itemize}

From this follows the division of an application into three libraries: the \emph{model},
which represents the game and its rules; the \emph{view}, which gathers the graphical
components; the \emph{controller}, \enquote{the conductor that links the two}.

\section{The first version}

In 2023 we carried out this project as a pair. Rather than dress up the library we were
given, we had chosen to write the whole game, to get a better grasp of GTK. The result was a
complete and pleasant game: three difficulty levels, random placement of the fleets, rules
and presentation screens, a polished visual identity.

But this game fitted in \emph{a single file} of \num{2666} lines, and its organisation
contradicted what the course taught:

\begin{itemize}
    \item the state of the game was spread across seven global variables and the graphical
          components themselves: the structure describing a cell held both
          \enquote{is there a ship?} and the address of the button that displays it;
    \item the three levels were three copies of one function several hundred lines long; a
          fix made in one was missing from the others;
    \item no rule could run without opening a window, so no automatic test was possible;
    \item the program only worked when started from its own folder, and its sound variant
          only on the machine where it had been written.
\end{itemize}

The program \emph{worked}. It could not be verified, nor modified without risk, nor reused.

\section{Goals of the rewrite}

Version~2 has a goal that is simple to state: \textbf{keep the game, redo everything else}.
The player must find the same rules, the same screens, the same texts. Beneath the surface,
the program must become what software maintained over time has to be:

\begin{enumerate}
    \item \textbf{modular} --- a model truly independent of the interface, to the point that
          a second interface can drive it;
    \item \textbf{verified} --- automatic tests for the rules, the computer's strategies and
          the interface itself;
    \item \textbf{reproducible} --- a game entirely determined by a seed, an identical build
          on a workstation and on a server;
    \item \textbf{deliverable} --- a single executable, packages, continuous integration and
          delivery.
\end{enumerate}

\section{Structure of the report}

\Cref{chap-exigences} sets out the rules of the game and the requirements.
\Cref{chap-architecture} presents the architecture and justifies it.
\Cref{chap-modele,chap-adversaire} detail the game engine and then the artificial opponent,
with the measurements of its strength. \Cref{chap-interface} describes the graphical
interface. \Cref{chap-qualite} sets out the verification approach and the continuous
integration pipeline. \Cref{chap-bilan} draws up the assessment in figures of the rewrite,
its limits and its possible extensions.
